\section{Method}
\label{sec:method}

A frozen language model learns $T$ tasks one after another, each with its own low-rank adapter, never
changed afterwards. A switch sends every input to one task (Figure~\ref{fig:overview}, Appendix~\ref{app:pictures}). We give three rules: a recogniser that
reads no instruction (Section~\ref{sec:recogmain}), a certified switch for inputs that carry one
(Section~\ref{sec:features}), and zones built from the inputs of a task (Section~\ref{sec:boundmain}).

\subsection{Setting: adapters, routes and answers}
\label{sec:setting}

Let $E$ be the input embedding table of the base model; $E$ and all other base weights $\theta_0$ are
frozen. Task $k$ adds a LoRA adapter \citep{hu2022lora} $\Delta W_k = \tfrac{\alpha}{r} B_k A_k$ to
every adapted weight matrix $W_0$. After $T$ tasks an adapted layer with input $z$ computes
\begin{align}
  h &\;=\; \Bigl( W_0 + \sum_{k=1}^{T} g_k(x)\, \Delta W_k \Bigr) z , \qquad g_k(x) \in \{0, 1\},
  \label{eq:layer}\\
  g_k(x) &\;=\; \mathbf{1}\bigl[\, k \le \rho_T(x) \,\bigr] , \label{eq:gates}\\
  y_T(x) &\;=\; F\bigl(x;\; \theta_0,\; \{A_k, B_k\}_{k \le \rho_T(x)}\bigr),
  \label{eq:answer}
\end{align}
where $F$ runs the model with greedy decoding. The methods differ only in the \emph{route}
$\rho_T(x) \in \{1, \dots, T\}$: an input with route $j$ is processed with the adapters of tasks
$1, \dots, j$, the model as it was right after task $j$ was learned.

\begin{lemma}[An answer changes only with its route]
\label{lem:route}
Assume that adapters are never changed after their task was learned, that the base model, its
tokenizer, the input template and truncation are fixed, and that decoding is deterministic. Then for
every rule that sets the route, every input $x$ and all $T' > T$: if $\rho_{T'}(x) = \rho_T(x)$, then
$y_{T'}(x) = y_T(x)$.
\end{lemma}

\begin{proof}
Let $j = \rho_T(x) = \rho_{T'}(x) \le T$. By Eq.~\eqref{eq:answer} both answers are
$F(x; \theta_0, \{A_k, B_k\}_{k \le j})$; these adapters are the same at both times, and $F$ is
deterministic.
\end{proof}

\subsection{Without instructions: a recogniser, an audit, and a handicap}
\label{sec:recogmain}

\paragraph{The price.} Let $\ell_i(T)$ be what a system has lost on task $i$ at time $T$ against the
best possible answers, and $p_i(x)$ the probability of input $x$ under task $i$. For two tasks
$j < k$ and every system whose answer depends on the input alone,
\begin{equation}
  \ell_j(T) + \ell_k(T) \;\ge\; \Omega_{jk} \;=\; \sum_x \min\bigl(p_j(x), p_k(x)\bigr)\, \gamma_{jk}(x),
  \label{eq:priceshort}
\end{equation}
where $\gamma_{jk}(x)$ is the least combined loss of one answer to $x$ under the two tasks: together
the two tasks lose at least the share of inputs that they have in common, weighted by how far their
right answers differ. This is zero when no input belongs to two tasks, as when every input begins
with its task's instruction and no instruction begins another
(Proposition~\ref{prop:noinstructionprice}, Appendix~\ref{sec:audit}).

\paragraph{The recogniser.} Let $u(x)$ be a frozen feature of the whole input (the mean of the
last-layer states of the frozen encoder of the base model, without any adapter, standardised with
the statistics of 2{,}000 instances of other SuperNI tasks, not $\mathcal{B}$) and $\bar u(x) = (u(x), 1)$. When task $k$ arrives, its training
inputs are split at random into fit inputs $D_k^{\mathrm{f}}$ and $m$ held-out inputs. The fit inputs
add one matrix and one vector to the store:
$G_k = \operatorname{mean}_{x \in D_k^{\mathrm{f}}} \bar u(x)\bar u(x)^{\!\top}$ and
$v_k = \operatorname{mean}_{x \in D_k^{\mathrm{f}}} \bar u(x)$, and after $T$ tasks, with a fixed $\varepsilon > 0$,
\begin{equation}
  V_T = \Bigl(\sum_{k \le T} G_k + \varepsilon I\Bigr)^{-1} [\,v_1, \dots, v_T\,], \qquad
  \rho_T(x) = \arg\max_{k \le T} \bigl(\bar u(x)^{\!\top} V_T\bigr)_k ,
  \label{eq:recogniserShort}
\end{equation}
the older task in a tie. $V_T$ is the ridge regression of the task labels on the features of all
tasks so far. The sums are only added, so no input is stored, $V_T$ does not depend on the order of
arrival, and every task has a column of its own and its adapter trains on all its inputs: no task
is \emph{blocked}. A recogniser of this kind is known from
class-incremental learning on frozen features \citep{mcdonnell2023ranpac,zhuang2022acil}; what we add is
what is said about its routes.

\paragraph{The audit.} For a task $j$ and two times $T < T'$, let $H$ be the number of its $m$
held-out inputs, which the rule never uses, whose route differs between $T$ and $T'$. Then $H/m$
estimates, for inputs drawn like the held-out ones, the share of changed routes, an upper limit for
changed answers (Lemma~\ref{lem:route}), and the Clopper--Pearson bound of $H$ in $m$ gives a
confidence bound for it (Proposition~\ref{prop:audit}). It needs later tasks that were not chosen by
looking at the held-out inputs, and it holds for one task and one pair of times. It can be read
before the newcomer's adapter is trained.

\paragraph{The handicap.} When task $k$ arrives, let $\xi_k(x)$ be the margin of the newcomer over
the best older task in $V_k$. Fix an integer $\nu \ge 0$. The threshold $\eta_k$ is the largest, over
older tasks $j$, of the $(\nu+1)$-th largest margin among the held-out inputs of task $j$, and never
below zero. The newcomer \emph{takes} an input if $\xi_k(x) > \eta_k$, and the route is the newest task
that took it (task 1 if none did). Both $V_k$ and $\eta_k$ are stored and never changed. The
newcomer pays: it may lose some of its own inputs. If later tasks are fixed without looking at them,
a new input of task $j$ drawn like its held-out inputs ever changes its route with probability at
most $(T_{\mathrm{end}}-j)(\nu+1)/(m+1)$ (Theorem~\ref{thm:handicap}, Appendix~\ref{sec:audit}).

These statements bound change, not quality: an input given to the wrong task from the start does not
change. The bound in advance is loose and grows with the number of later tasks (on the 15
development tasks, \AudDevHPrimeTTwoBound{}\% against \AudDevHPrimeTTwoChanged{}\% measured), and the audit is not valid
for the same held-out inputs that set the thresholds.

\subsection{With instructions: protected zones and the certified switch}
\label{sec:features}\label{sec:zones}\label{sec:switch}\label{sec:guarantee}

When every input carries the instruction of its task, a second rule gives an exact statement. In
instruction tuning every input of task $k$ is the task's instruction followed by the instance; the
part before the instance is the \emph{opening} $o_k$, of $n_k$ tokens (Appendix~\ref{app:impl}), and
an input that begins with it carries the \emph{original wording}, or is \emph{unedited}. The rule
reads an input through the frozen embedding table: $\psi_n(x)$ is the weighted average
\citep{arora2017simple} of the embeddings of its first $n$ tokens, rescaled with the mean and
covariance of a fixed collection $\mathcal{B}$ of unrelated inputs (Eqs.~\eqref{eq:weights}
to~\eqref{eq:whiten}, Appendix~\ref{app:method}). It depends only on the base model and on
$\mathcal{B}$, and never changes.

\begin{definition}[Signature, radius, zone]
When task $k$ is learned, it stores its \emph{signature} $s_k = \psi_{n_k}(o_k)$, which equals
$\psi_{n_k}(x)$ for every unedited input $x$ of task $k$. Let $\mathcal{S}_k$ be the signatures known
then: those of the unrelated tasks that $\mathcal{B}$ is drawn from, and $s_1, \dots, s_{k-1}$. The
\emph{radius}, the \emph{protected zone} and the \emph{relative distance} are
\begin{align}
  s_k &= \psi_{n_k}(o_k), \label{eq:signature}\\
  R_k &\;=\; \kappa \cdot \min_{s \in \mathcal{S}_k} \lVert s_k - s \rVert_2 , \quad \kappa = \tfrac{1}{2}, \label{eq:radius}\\
  Z_k &= \{\, x : \delta_k(x) \le 1 \,\} , \qquad
  \delta_k(x) \;=\; \lVert \psi_{n_k}(x) - s_k \rVert_2 \,/\, R_k . \label{eq:zone}
\end{align}
If $R_k = 0$, then $\delta_k(x) = 0$ at distance zero and $\infty$ elsewhere.
\end{definition}

Only $(n_k, s_k, R_k)$ is kept per task; the zone reaches half-way to the nearest task known.

\begin{definition}[Route and certificate]
After $T$ tasks, the \emph{certificate} of an input and its \emph{route} are
\begin{align}
  c_T(x) &\;=\; \mathbf{1}\bigl[\, x \in Z_1 \cup \dots \cup Z_T \,\bigr] , \label{eq:certificate}\\[2pt]
  \rho_T(x) &\;=\;
  \begin{cases}
    \min \{\, k \le T : x \in Z_k \,\} & \text{if } c_T(x) = 1 \quad \text{(the oldest zone that contains } x \text{)}, \\[2pt]
    \arg\min_{k \le T}\; \delta_k(x) & \text{if } c_T(x) = 0 \quad \text{(the nearest zone)} .
  \end{cases}
  \label{eq:route}
\end{align}
Ties go to the older task.
\end{definition}

An input inside a zone is \emph{certified}; one outside every zone goes to the nearest zone, without a
certificate. This rule is \cert{}. When task $t$ is learned, its training inputs inside an older zone
are \emph{blocked}: they go to the older task and give no gradient to the new adapter; their number
is $b_t$. Otherwise training is that of plain adapters. For comparison, \plain{}
switches all adapters on; \keyswitch{} goes to the task with the longest opening that matches token
for token (all on if none); \nearest{} uses $\arg\min_k \delta_k(x)$ for every input; \centre{} picks
the nearest centre, the centre of a task being the mean word-vector feature of its whole training
inputs
(Algorithm~\ref{alg:switch} and Table~\ref{tab:methods}, Appendix~\ref{app:method}).

\begin{theorem}[Write-once answers]
\label{thm:writeonce}
Assume that zones and adapters are never changed once created, that the base model, its tokenizer,
the input template and truncation, and the quantities above are fixed, and that decoding is
deterministic. Let $c_T(x) = 1$. Then for every later time $T' \ge T$ and every choice of tasks
$T+1, \dots, T'$: $c_{T'}(x) = 1$, $\rho_{T'}(x) = \rho_T(x)$ and $y_{T'}(x) = y_T(x)$.
\end{theorem}

\begin{proof}
Let $j = \rho_T(x)$. The zones $Z_1, \dots, Z_T$ are the same sets at time $T'$, so
$\{k \le T' : x \in Z_k\}$ is the union of $\{k \le T : x \in Z_k\}$, non-empty with minimum $j$, and
indices larger than $T$. Its minimum is $j$, so $c_{T'}(x) = 1$ and $\rho_{T'}(x) = j$, and
Lemma~\ref{lem:route} gives $y_{T'}(x) = y_T(x)$.
\end{proof}

The proof uses only determinism, nothing else about the model or later tasks (floating point:
Appendix~\ref{app:method}). Four more statements are proved in Appendix~\ref{app:method}. Every input
that begins with the opening of task $k$ lies in $Z_k$ (Proposition~\ref{prop:own}). An edit that
removes or adds at most the share $R_k/D$ of the weight of the opening ($D$: the largest distance between rescaled
embedding rows) keeps an input in its zone, a worst-case bound (Proposition~\ref{prop:edits}). The training inputs of a new task inside an older
zone are known before it is trained, and if there are none, training is plain training
(Proposition~\ref{prop:price}). An input outside every zone changes its route only if the newest zone
is nearer to it than all older ones (Proposition~\ref{prop:uncertified}). A zone variant for tasks with
nearly identical instructions, \exactfirst{}, is in Appendix~\ref{sec:exactfirst}.

\subsection{Zones from the inputs of a task}
\label{sec:boundmain}

A third rule builds each zone from the inputs of its task, calibrated on held-out inputs, and needs no
instruction. Its bound does not grow with the number of later tasks (Theorem~\ref{thm:bound},
Appendix~\ref{sec:bound}). In our tests it held, but with frozen word-vector features the zones of
early tasks held inputs of later tasks, which were then blocked.
